\documentclass[11pt]{article}

\usepackage[margin=1in]{geometry}
\usepackage{booktabs}
\usepackage{tabularx}
\usepackage{array}
\usepackage{parskip}
\usepackage{hyperref}
\usepackage{enumitem}

\hypersetup{
  colorlinks=true,
  linkcolor=black,
  urlcolor=blue
}

\newcolumntype{L}[1]{>{\raggedright\arraybackslash}p{#1}}
\newcolumntype{Y}{>{\raggedright\arraybackslash}X}

\title{\textbf{VPRTempo Experiment Comparison Report}\\[4pt]
Transfer Evaluation on GardensPoint and Campus Datasets}
\author{Samuel Olusola}
\date{April 24, 2026}

\begin{document}
\maketitle

\section*{Objective}
This report summarizes the first \textbf{VPRTempo} experiments run inside this repository and compares them against the descriptor runs already saved in \texttt{output\_images/runs}. The goal was to answer a practical question: \textbf{does the official VPRTempo pretrained model transfer well enough to be competitive with the strongest descriptors we have already tested?}

\section*{Important Caveat}
This comparison must be interpreted carefully. The VPRTempo package provides a \textbf{dataset-specific demo checkpoint}, \texttt{springfall\_VPRTempo\_IN3136\_FN6272\_DB500.pth}, which is trained on a \textbf{Nordland seasonal place-recognition setup}. It is not a universal pretrained global descriptor in the same sense as SALAD or SelaVPR++.

For integration with this repository, we used the hidden \textbf{feature-layer activations} as the descriptor exported into the cosine-similarity pipeline. That makes the experiment a \textbf{transfer probe}, not a faithful reproduction of the original VPRTempo evaluation protocol.

\section*{What Was Run}
The following runs were completed with the official demo checkpoint:
\begin{itemize}[leftmargin=*]
  \item GardensPoint benchmark: \texttt{output\_images/runs/benchmarks/0012\_demo\_gardenspoint\_vprtempo\_20260424\_213108}
  \item Campus strict image-level evaluation: \texttt{output\_images/runs/campus\_strict/0007\_campus\_strict\_vprtempo\_20260424\_213108}
  \item Campus place-level evaluation: \texttt{output\_images/runs/campus\_place\_level/0006\_campus\_place\_level\_vprtempo\_20260424\_213108}
\end{itemize}

The comparison below uses the other result files already available under the corresponding \texttt{runs/} folders.

\section*{Benchmark Comparison: GardensPoint}
\begin{center}
\begin{tabularx}{\textwidth}{L{2.8cm} Y c c c c}
\toprule
\textbf{Descriptor} & \textbf{Run} & \textbf{AUC} & \textbf{R@1} & \textbf{R@5} & \textbf{R@10} \\
\midrule
CosPlace & \texttt{0001\_...cosplace} & 0.592 & 0.350 & 0.840 & 0.925 \\
EigenPlaces & \texttt{0002\_...eigenplaces} & 0.675 & 0.365 & 0.840 & 0.940 \\
NetVLAD & \texttt{0005\_...netvlad} & 0.199 & 0.185 & 0.480 & 0.695 \\
PatchNetVLAD & \texttt{0006\_...patchnetvlad} & 0.802 & 0.440 & 0.800 & 0.840 \\
AlexNet & \texttt{0007\_...alexnet} & 0.195 & 0.300 & 0.665 & 0.745 \\
HDC-DELF & \texttt{0008\_...hdc\_delf} & 0.742 & 0.420 & 0.870 & 0.925 \\
SAD & \texttt{0009\_...sad} & 0.030 & 0.090 & 0.225 & 0.305 \\
SALAD & \texttt{0010\_...salad} & 0.924 & 0.450 & 0.955 & 0.990 \\
SelaVPR++ & \texttt{0011\_...selavpr} & \textbf{0.937} & \textbf{0.465} & 0.950 & \textbf{0.995} \\
VPRTempo & \texttt{0012\_...vprtempo} & 0.065 & 0.095 & 0.230 & 0.290 \\
\bottomrule
\end{tabularx}
\end{center}

\textbf{Observation.} VPRTempo is \textbf{not competitive} on GardensPoint in this transfer setting. Its performance is only slightly above SAD and far below every strong modern descriptor tested in the repository.

\section*{Campus Comparison: Strict Image-Level}
\begin{center}
\begin{tabularx}{\textwidth}{L{2.8cm} c c c c c}
\toprule
\textbf{Descriptor} & \textbf{AUC} & \textbf{R@100P} & \textbf{R@1} & \textbf{R@5} & \textbf{R@10} \\
\midrule
CosPlace & 0.470 & 0.030 & 0.727 & 1.000 & 1.000 \\
PatchNetVLAD & 0.420 & 0.000 & 0.727 & 0.970 & 0.970 \\
HDC-DELF & 0.488 & 0.030 & 0.788 & 0.939 & 0.970 \\
SALAD & 0.504 & 0.030 & 0.818 & 1.000 & 1.000 \\
SelaVPR++ & 0.464 & \textbf{0.061} & \textbf{0.909} & 1.000 & 1.000 \\
VPRTempo & 0.064 & 0.000 & 0.182 & 0.515 & 0.606 \\
\bottomrule
\end{tabularx}
\end{center}

\textbf{Observation.} Under the strict 1-to-1 protocol, VPRTempo again performs very poorly. It fails both as a high-quality retriever and as a no-match rejector. In contrast, the stronger modern descriptors preserve very high top-\(k\) retrieval even under this harsh setup.

\section*{Campus Comparison: Place-Level}
\begin{center}
\begin{tabularx}{\textwidth}{L{2.8cm} c c c c c}
\toprule
\textbf{Descriptor} & \textbf{AUC} & \textbf{R@100P} & \textbf{R@1} & \textbf{R@5} & \textbf{R@10} \\
\midrule
HDC-DELF & 0.507 & 0.132 & 0.922 & 0.969 & 0.984 \\
CosPlace & 0.627 & 0.114 & 0.969 & 1.000 & 1.000 \\
PatchNetVLAD & 0.641 & 0.074 & 0.938 & 1.000 & 1.000 \\
SALAD & 0.653 & 0.130 & 0.969 & 1.000 & 1.000 \\
SelaVPR++ & \textbf{0.658} & \textbf{0.158} & 0.969 & 1.000 & 1.000 \\
VPRTempo & 0.151 & 0.000 & 0.219 & 0.641 & 0.734 \\
\bottomrule
\end{tabularx}
\end{center}

\textbf{Observation.} Even after relabeling credits multiple valid views of the same place, VPRTempo still transfers poorly. This is important because it shows the weak result is not mainly due to the old strict evaluation protocol.

\section*{Why the Models Behave Differently}
\subsection*{1. SALAD and SelaVPR++ perform best because they are strong general VPR descriptors}
These two models are the strongest overall because they were designed as \textbf{modern global place-recognition descriptors}. They use strong pretrained visual backbones and learn representations that remain useful across datasets. In practice:
\begin{itemize}[leftmargin=*]
  \item SALAD gives the strongest overall \textbf{controlled benchmark transfer} among the earlier models we tested.
  \item SelaVPR++ slightly surpasses SALAD on GardensPoint and is also very strong on both campus protocols.
  \item Their representations are much closer to the type of \textbf{drop-in descriptor} our current cosine-similarity pipeline expects.
\end{itemize}

\subsection*{2. CosPlace remains a very strong lightweight baseline}
CosPlace is older than SALAD and SelaVPR++, but it remains highly competitive, especially on the relabeled campus dataset. This is a good sign that the campus place-level task is still structurally a standard VPR retrieval problem where robust learned global descriptors work well.

\subsection*{3. PatchNetVLAD gives high true-positive counts but also more unstable threshold behavior}
PatchNetVLAD often finds correct places, but its thresholded behavior is less clean, especially on campus. That is consistent with local-patch matching being more sensitive to repeated structures, local clutter, and viewpoint noise.

\subsection*{4. HDC-DELF is respectable but no longer state-of-the-art in this setup}
HDC-DELF remains surprisingly solid, especially on place-level campus evaluation, but it does not reach the same combination of benchmark strength and clean transfer shown by SALAD and SelaVPR++.

\subsection*{5. VPRTempo performs poorly because this is the wrong checkpoint for this use case}
The main reason for VPRTempo's weak result is not that the method is ``bad'' in general. It is that the experiment is a \textbf{cross-domain transfer test using a dataset-specific spiking model checkpoint}.

There are several reasons this hurts it badly here:
\begin{itemize}[leftmargin=*]
  \item The checkpoint is trained for a \textbf{Nordland seasonal route}, not for GardensPoint or campus scenes.
  \item VPRTempo is fundamentally closer to a \textbf{task-specific place recognizer} than a universal descriptor.
  \item Its preprocessing pipeline is very specialized: grayscale conversion, gamma correction, patch normalization, and spike encoding.
  \item Our integration exports the hidden \textbf{feature-layer activations} into a standard cosine-retrieval pipeline, which is a useful probe but not the model's native end-to-end evaluation mode.
\end{itemize}

So the weak result is \textbf{expected and explainable}. It does not invalidate the VPRTempo paper; it simply shows that the official demo checkpoint does not transfer as a strong generic descriptor in our current pipeline.

\section*{Main Conclusions}
\begin{enumerate}[leftmargin=*]
  \item \textbf{VPRTempo should not be treated as a direct drop-in replacement} for SALAD or SelaVPR++ in this repository.
  \item \textbf{SALAD and SelaVPR++ are currently the strongest descriptors} we have tested for our day-to-night and campus transfer experiments.
  \item \textbf{CosPlace remains a strong practical baseline}, especially when the evaluation better reflects real place-level matching.
  \item The VPRTempo result is still useful because it answers an important question: \textbf{the official demo checkpoint does not transfer well to our current datasets}.
\end{enumerate}

\section*{Recommendation}
If the goal is to improve the current project quickly, the next priority should stay with \textbf{SALAD, SelaVPR++, and related descriptor-family experiments}. VPRTempo should only be revisited if we are willing to treat it as a \textbf{separate branch} that may require:
\begin{itemize}[leftmargin=*]
  \item dataset-specific training,
  \item its own evaluation protocol,
  \item and a robotics/CPU-efficiency motivation rather than a universal descriptor motivation.
\end{itemize}

\section*{Files Used}
\begin{itemize}[leftmargin=*]
  \item VPRTempo run results: \texttt{output\_images/runs/benchmarks/0012\_demo\_gardenspoint\_vprtempo\_20260424\_213108/results.txt}
  \item Strict campus VPRTempo: \texttt{output\_images/runs/campus\_strict/0007\_campus\_strict\_vprtempo\_20260424\_213108/results.txt}
  \item Place-level campus VPRTempo: \texttt{output\_images/runs/campus\_place\_level/0006\_campus\_place\_level\_vprtempo\_20260424\_213108/results.txt}
\end{itemize}

\end{document}
